\documentclass[10pt,twocolumn]{article}
\usepackage[letterpaper,margin=0.65in,columnsep=0.3in]{geometry}
\usepackage{amsmath,amssymb}
\usepackage{times}
\usepackage{xcolor}
\usepackage[compact]{titlesec}
\usepackage{enumitem}
\usepackage[hidelinks]{hyperref}
\definecolor{accent}{RGB}{30,80,140}
\titleformat{\section}{\normalsize\bfseries\color{accent}}{\thesection}{0.6em}{}
\titlespacing*{\section}{0pt}{8pt}{3pt}
\setlist{nosep,leftmargin=1em,itemsep=1.5pt}
\setlength{\parindent}{0pt}
\setlength{\parskip}{4pt}
\setlength{\abovedisplayskip}{3pt}
\setlength{\belowdisplayskip}{3pt}
\pagestyle{empty}

\begin{document}

\twocolumn[{%
\begin{center}
{\Large\bfseries Probabilistic Graphical Models of\\[1pt] U.S. Senate Roll-Call Voting}\\[6pt]
{\small CSE 150A/250A Project Proposal \quad$\cdot$\quad \textbf{Group:} \textit{[Group Name]}}\\[1pt]
{\small \textit{[Name 1 (PID)], [Name 2 (PID)], [Name 3 (PID)], [Name 4 (PID)], [Name 5 (PID)]}}
\end{center}
\vspace{6pt}
}]

\section{Problem Description}
Senators do not vote independently: party discipline, ideological blocs and cross-party alliances
create strong dependencies among votes, and some votes are missing. We ask three questions:
\begin{itemize}
\item \textbf{Q1:} given the observed votes of some senators on a roll call, what is
$P(X_i{=}\text{Yea}\mid\mathbf{x}_O)$ for each remaining senator?
\item \textbf{Q2:} which dependency structure (latent classes, a tree, or a loopy graph) best explains the data?
\item \textbf{Q3:} do the learned dependencies carry over from one Congress to the next?
\end{itemize}
Answering them requires explicit independence assumptions, learning with and without missing data
(MLE vs.\ EM), and exact vs.\ approximate inference. Our aim is to understand what each model captures
or misses, not to maximize accuracy.

\section{Dataset Source}
\textbf{Voteview} roll-call data for the 117th (2021--22) and 118th (2023--24) Senates
(Lewis et al., \url{https://voteview.com/data}): individual votes, member party/state, and roll-call metadata.
\begin{itemize}
\item \textbf{Variables:} senator $X_i\in\{-1,+1\}$ (Nay/Yea), $N\approx100$; each roll call is one sample:
949 (117th) and 691 (118th) roll calls; 91 senators served in both.
\item \textbf{Preprocessing:} cast codes 1--3$\to$Yea, 4--6$\to$Nay, 7--9$\to$missing (4--5\%, assumed
missing at random); drop near-unanimous votes ($\approx$4\% of roll calls); drop senators present for
under half of a Congress; chronological 80/20 split
within each Congress. Party and state are used only to interpret results.
\item \textbf{Limitations:} few samples relative to $O(N^2)$ pairwise parameters; strategic absences
make missingness not truly random; membership changes between Congresses.
\end{itemize}

\section{Methodology}
We compare three models of increasing expressiveness.

\textbf{M1: Latent-class Naive Bayes.} Each roll call has a hidden coalition type $Z\in\{1,\dots,K\}$
(e.g.\ Democratic-backed, Republican-backed, bipartisan), and votes are independent given $Z$:
\[
P(z,\mathbf{x})=\pi_z\textstyle\prod_i P(x_i\mid z).
\]
Parameters $\pi_k$ and $\theta_{ik}=P(X_i{=}{+}1\mid Z{=}k)$ are learned by \textbf{EM}. The E-step
computes responsibilities $\gamma_{mk}=P(Z^{(m)}{=}k\mid\mathbf{x}_O^{(m)})$ from observed votes only, so
missing data is handled directly. The M-step sets $\pi_k$ to the average $\gamma_{mk}$ and $\theta_{ik}$ to the
$\gamma$-weighted Yea rate. We use several random restarts and choose $K\in\{2,\dots,6\}$ by held-out
log-likelihood. Exact inference costs $O(KN)$.

\textbf{M2: Chow-Liu tree BN.} A maximum-weight spanning tree on empirical pairwise mutual information
$I(X_i;X_j)$, rooted into a DAG; this is the maximum-likelihood tree. CPTs are fit by \textbf{MLE} on
complete data and by \textbf{EM} with missing votes. Exact marginals come from sum-product belief
propagation in $O(N)$.

\textbf{M3: Sparse Ising model.} Trees cannot represent loops of allied senators, so we fit
\[
P(\mathbf{x})\propto\exp\Big(\textstyle\sum_i\theta_i x_i+\sum_{i<j}W_{ij}x_ix_j\Big),
\]
where $W_{ij}>0$ ($<0$) means senators $i$ and $j$ tend to agree (disagree). The partition function sums
over $2^N$ vote patterns, so we learn $(\theta,W)$ by $\ell_1$-regularized \textbf{pseudo-likelihood}, which
reduces to one convex logistic regression per senator; $\lambda$ is chosen on held-out data. Inference
uses \textbf{Gibbs sampling} and naive \textbf{mean-field}, both checked against \textbf{exact
enumeration} on a 16-senator model.

\section{Experiments}
\begin{itemize}
\item \textbf{E1, exact vs.\ approximate:} on the 16-senator model, scale couplings by
$\beta\in[0.1,2]$; report $\ell_1$ and KL error of Gibbs and mean-field marginals vs.\ runtime.
\item \textbf{E2, missing data:} mask 10/30/50\% of training votes; compare complete-case MLE and EM
(M1, M2) by held-out log-likelihood and imputation accuracy.
\item \textbf{E3, model comparison:} reveal a random 20\% of votes on test roll calls and predict the
rest; report accuracy and per-vote NLL against a per-senator baseline, and interpret learned classes,
tree edges and couplings using party and state.
\item \textbf{E4, cross-Congress transfer:} train on the 117th Senate and test on the 118th, restricted to
senators who served in both. We compare against training on the 118th itself, and test which parameters
transfer: for M1, refit only the coalition mix $\pi$ while keeping $\theta$; for M1 and M2, use 117th CPTs
as Beta priors (pseudo-counts) and measure how many 118th roll calls are needed to match in-Congress
performance.
\end{itemize}
\textbf{Stretch goal:} stochastic EM for M3 (Gibbs imputation of missing votes alternating with
pseudo-likelihood refitting), added to E2.

\end{document}
